\documentclass[10pt,a4paper]{article}
\usepackage{amsmath,amssymb,amsthm,graphicx}
\usepackage{hyperref}
\usepackage{latexsym}
\usepackage{amssymb}
\usepackage{amsthm}
\usepackage{ulem}
\usepackage{enumitem}
\setlist{itemsep=0.4pt,topsep=0pt}
\parindent 0pt
\parindent 0pt
\parskip 12pt
\newcommand{\HCF}{\mathrm{hcf}}
\newcommand{\ZZ}{\mathbb{Z}}
\newcommand{\rk}{\mathrm{rk}}
\DeclareMathOperator{\len}{length}
\DeclareMathOperator{\bl}{b}
\DeclareMathOperator{\cl}{cl}
\DeclareMathOperator{\di}{d}
\DeclareMathOperator{\Cl}{Cl}
\DeclareMathOperator{\acl}{acl}
\DeclareMathOperator{\tp}{tp}
\DeclareMathOperator{\Th}{Th}
\newcommand{\Cor}{\mathcal{C}_\mathrm{or}}
\newcommand{\Cun}{\mathcal{C}_\mathrm{un}}
\newcommand{\mons}{\mathbb{M}}
\newcommand{\strong}{\mathrm{strong}}
\newcommand{\axioms}{\Delta_{\mathrm{FSTS}}}
\newtheoremstyle{note}% hnamei
{8pt}% hSpace abovei
{14pt}% hSpace belowi
{\normalfont}% hBody fonti
{}% hIndent amounti1
{\bfseries}% hTheorem head fonti
{ }% hPunctuation after theorem headi
{.5em}% hSpace after theorem headi2
{}% hTheorem head spec (can be left empty, meaning ‘normal’)i
\theoremstyle{note}
\newtheorem{thm}{Theorem}[section]
%\newtheorem{theorem}{Theorem}
\newtheorem{lemma}[thm]{Lemma} 
\newtheorem{df}[thm]{Definition}
\newtheorem{rmk}[thm]{Remark} 
\newtheorem{prop}[thm]{Proposition}
\newtheorem{fact}[thm]{Fact}
\newtheorem{cor}[thm]{Corollary}
\newtheorem{conj}[thm]{Conjecture}
\newtheorem{ex}[thm]{Example}
\newtheorem{qu}[thm]{Question}
\newtheorem*{notation}{Notation} 
\title{An HF-ordering over a partial subsystem exists exactly when that subsystem can be made an initial segment}
\author{Silvia Barbina, Enrique Casanovas}
\date{arxiv:2411.13723v3; Creative Commons Attribution 4.0 International}
\begin{document}
\maketitle
\noindent\textit{Editorial note.} The counted statement is the author's labelled lemma environment at source0.tex lines 333-335; its proof environment opens at line 336, immediately after \textbackslash{}end\{lemma\}, so the pairing is by adjacency and the retained proof range 336-339 is the whole environment.\par\medskip
\noindent\textit{Editorial note.} This source is newly approved for this line; the preamble is built from the author's own package lines, macro block and theorem-environment declarations (source lines 4-65), with the author's optioned amsmath line omitted so that it cannot clash with the wrapper's amsmath.\par\medskip
\noindent\textit{Editorial note.} The source carries no label-number dump and the retained counter values are left empty, so the printed numbering of the card is the wrapper's default; the identity of the counted item is its author label order-over.\par\medskip
\noindent\textit{Editorial note.} Metadata provenance: the source record of the pinned copy and the licence evidence file are both present; the licence is CC BY 4.0 with status supported.\par\medskip
\noindent\textit{Editorial note.} r002 note: r001 failed because the retained author block carried the layout setting mathindent, which only exists when amsmath is loaded with the fleqn option; r002 starts that span at line 22 instead and changes nothing else.\par\medskip
\noindent\textit{Editorial note.} r003 note: the pre-publication gate compares the card's source\_version with the source\_version recorded in the licence evidence of the same source. This source's evidence file records exactly "2411.13723v3", while the card had carried "arXiv:2411.13723v3 (CC BY 4.0)" copied from the source record; r003 aligns the card with the evidence and changes nothing else. The evidence file itself is not modified.\par\medskip
\noindent\textit{Editorial scope.} One original result of the pinned author manuscript, counted once: the equivalence between the existence of an HF-ordering over a partial subsystem and the possibility of making that subsystem an initial segment (author label order-over).\par\medskip
\noindent\textit{Allowed standard prerequisites (editorial).}\begin{itemize}
\item Partial Steiner triple system. A partial Steiner triple system is a structure with a partial binary operation whose defined products are the blocks, with the Steiner condition on the operation.. It is the structure that carries the orderings.
\item HF-ordering. An HF-ordering of B over A is a linear ordering in which the closure of A is A itself and every element is either in A or produces a block whose other members are smaller.. It is the notion the lemma compares with having A as an initial segment.
\item Initial segment. A subset is an initial segment of an ordering when every element below an element of the subset belongs to the subset.. It is the second description of the ordering.
\end{itemize}\medskip
\begin{lemma}\label{order-over}
	Let $A\subseteq B$ be partial STSs. There is an HF-ordering of $B$ over $A$ if and only if there is an HF-ordering of $B$ where $A$ is an initial segment. In fact, given any HF-ordering $<$ of $B$ over $A$, the sum $(A,<\restriction A)+(B\smallsetminus A, <\restriction (B\smallsetminus A))$ is an HF-ordering of $B$ with initial segment $A$.
\end{lemma}	

\begin{proof}
	On the one hand, if $<$ is an HF-ordering of $B$ with initial segment $A$, it easy to see that it is an HF-ordering over $A$. On the other hand, assume $<$ is an HF-ordering of $B$ over $A$. Then the only way in which the sum can fail to be an HF-ordering is for some $a\in B \smallsetminus A$ to be a product of two different pairs of smaller elements, %each in $B\smallsetminus A$ and smaller than $a$ in $<$, or in $A$, contradicting that $<$ is an HF-ordering over $A$.
and each element is either smaller than $a$ in $<$ or belongs to $A$, contradicting that $<$ is an HF-ordering over~$A$.
\end{proof}

\end{document}
